\documentclass[11pt,a4paper]{article}
\usepackage[margin=25mm,headheight=14pt]{geometry}
\usepackage{fontspec}
\setmainfont{lmroman10-regular.otf}[BoldFont=lmroman10-bold.otf,ItalicFont=lmroman10-italic.otf,BoldItalicFont=lmroman10-bolditalic.otf]
\setsansfont{lmsans10-regular.otf}
\setmonofont{lmmono10-regular.otf}[Scale=MatchLowercase]
\newfontfamily\authorsymbol{Noto Sans Symbols}[Scale=MatchUppercase]
\usepackage{amsmath,amssymb,amsthm,mathtools,booktabs}
\usepackage{microtype}
\usepackage[hidelinks,unicode]{hyperref}
\usepackage{fancyhdr}
\pagestyle{fancy}\fancyhf{}\fancyhead[L]{\small\itshape Finite peak certificates for damped recurrences}\fancyhead[R]{\small\thepage}\renewcommand{\headrulewidth}{0pt}
\setlength{\parindent}{1.2em}\setlength{\parskip}{0pt}\setlength{\emergencystretch}{2em}
\newtheorem{theorem}{Theorem}\newtheorem{proposition}[theorem]{Proposition}\newtheorem{lemma}[theorem]{Lemma}\newtheorem{corollary}[theorem]{Corollary}\theoremstyle{definition}\newtheorem{definition}[theorem]{Definition}\newtheorem{example}[theorem]{Example}
\DeclareMathOperator{\rank}{rank}\DeclareMathOperator{\spanop}{span}\DeclareMathOperator{\diag}{diag}
\newcommand{\R}{\mathbb R}\newcommand{\C}{\mathbb C}\newcommand{\Z}{\mathbb Z}\newcommand{\N}{\mathbb N}\newcommand{\ip}[2]{\langle #1,#2\rangle}\newcommand{\norm}[1]{\lVert #1\rVert}\newcommand{\one}{\mathbf 1}
\hypersetup{pdftitle={Sharp Two-Sample Bounds and Finite Peak Certificates for Damped Recurrences},pdfauthor={lmtlssss🜐},pdfsubject={Mathematics: linear recurrences and exact certification},pdfkeywords={linear recurrence, sharp bound, observability, peak certificate, Stein equation}}
\title{\LARGE Sharp Two-Sample Bounds and Finite\\Peak Certificates for Damped Recurrences}
\author{lmtlssss{\authorsymbol 🜐}}
\date{2 November 2026\\[3pt]\small\href{https://doi.org/10.5281/zenodo.23897133}{10.5281/zenodo.23897133}}
\begin{document}\maketitle\thispagestyle{plain}

\begin{abstract}
Consider a real second-order recurrence with distinct stable conjugate roots
$a\pm ib$, where $b>0$ and $a^2+b^2<1$.
We prove the sharp two-sample bound
$\max\{x_0^2,x_1^2\}\ge
 b^2R^2/(1+a^2+b^2+2|a|)$,
where $R^2=x_0^2+((x_1-a x_0)/b)^2$.
It yields a finite, coefficient-dependent prefix containing the absolute
maximum of every infinite trajectory. An adaptive rational certificate can
terminate earlier by comparing the observed peak with the remaining
geometric envelope. We extend the sharp observation constant to arbitrary
finite sample sets through a two-dimensional strip polygon, and combine the
peak certificate with exact finite quadratic-sum formulas.
A worked rational recurrence attains its infinite absolute peak $6/5$ at
index two; four samples certify the peak, and its infinite square sum is
$175/39$. Exact independent computations accompany the proofs.
\end{abstract}
\noindent\textit{Keywords.} Linear recurrence; damped oscillation; sharp
observability bound; finite peak certificate; quadratic sums; rational arithmetic.

\section{A two-dimensional recurrence representation}
Fix $a\in\R$ and $b>0$ such that
\begin{equation}\label{eq:stable}
 r^2=a^2+b^2<1.
\end{equation}
Thus $0<r<1$. Consider arbitrary real initial values $x_0,x_1$ and
\begin{equation}\label{eq:recurrence}
 x_{n+2}=2a x_{n+1}-r^2x_n\qquad(n\ge0).
\end{equation}
Set
\begin{equation}\label{eq:state}
 z=\begin{pmatrix}u\\v\end{pmatrix}
 =\begin{pmatrix}x_0\\(x_1-a x_0)/b\end{pmatrix},\quad
 S=\begin{pmatrix}a&b\\-b&a\end{pmatrix},\quad
 R^2=u^2+v^2.
\end{equation}
Then $x_n=e_1^\mathsf TS^n z$. The characteristic equation of $S$ gives
\eqref{eq:recurrence}, and its first two outputs are $x_0,x_1$.
Since $S^\mathsf TS=r^2I$,
\begin{equation}\label{eq:envelope}
 |x_n|\le Rr^n.
\end{equation}
Write $(a+ib)^k=a_k+ib_k$. Equivalently,
$x_k=a_ku+b_kv$.

\section{The sharp two-sample constant}
\begin{theorem}[Sharp adjacent-sample bound]\label{thm:two}
For all real $x_0,x_1$,
\begin{equation}\label{eq:sharp}
 \boxed{\max\{x_0^2,x_1^2\}\ge c^2R^2,
 \qquad c^2=\frac{b^2}{1+r^2+2|a|}.}
\end{equation}
The coefficient $c^2$ is the largest constant for which this inequality
holds for every pair of initial values.
\end{theorem}
\begin{proof}
Put $M=\max\{|x_0|,|x_1|\}$. If $M=0$, the result is immediate.
For fixed $M$, the squared radius is
\[
 R^2=x_0^2+\frac{(x_1-a x_0)^2}{b^2}.
\]
This convex quadratic function attains its maximum over the square
$[-M,M]^2$ at a vertex. Choosing the endpoint signs so that
$|x_1-a x_0|=(1+|a|)M$ gives
\[
 R^2\le M^2\left(1+\frac{(1+|a|)^2}{b^2}\right)
 =M^2\frac{1+r^2+2|a|}{b^2}.
\]
Rearrangement proves~\eqref{eq:sharp}. Equality is attained with
$x_0=M$ and $x_1=-\operatorname{sgn}(a)M$ when $a\ne0$; for $a=0$,
any endpoint signs attain equality. Hence the constant is sharp.
\end{proof}

\begin{corollary}[Separated samples]\label{cor:separated}
Let $k\ge1$ and $b_k\ne0$. Then
\begin{equation}\label{eq:separated}
 \max\{x_0^2,x_k^2\}\ge c_k^2R^2,\qquad
 c_k^2=\frac{b_k^2}{1+r^{2k}+2|a_k|}.
\end{equation}
The constant is sharp. If $b_k=0$, these two samples alone admit no
positive constant multiplying $R^2$.
\end{corollary}
\begin{proof}
The pair $(x_0,x_k)=(u,a_ku+b_kv)$ has the same form used in
Theorem~\ref{thm:two}; replace $a,b$ there by $a_k,b_k$ and use
$a_k^2+b_k^2=r^{2k}$. If $b_k=0$, choose $u=0$ and $v\ne0$.
Both observed samples vanish while $R>0$.
\end{proof}

\section{A finite prefix for the infinite absolute peak}
For a nonzero trajectory, its absolute peak is positive and attained,
because~\eqref{eq:envelope} tends to zero.

\begin{theorem}[Uniform prefix certificate]\label{thm:prefix}
Let $N\ge2$ satisfy $r^{2N}<c^2$, with $c^2$ as in~\eqref{eq:sharp}.
Then every solution of~\eqref{eq:recurrence} satisfies
\begin{equation}\label{eq:prefix}
 \boxed{\max_{n\ge0}|x_n|=\max_{0\le n<N}|x_n|.}
\end{equation}
One may choose
\begin{equation}\label{eq:N}
 N=\max\left\{2,1+\left\lfloor
             \frac{\log c^2}{\log r^2}\right\rfloor\right\}.
\end{equation}
For separated samples, the same assertion holds whenever $N>k$ and
$r^{2N}<c_k^2$.
\end{theorem}
\begin{proof}
If $R=0$, all values vanish. Otherwise, for every $n\ge N$,
\[
 |x_n|^2\le R^2r^{2n}\le R^2r^{2N}<R^2c^2
 \le\max\{x_0^2,x_1^2\}.
\]
Both comparison samples belong to the prefix. This proves~\eqref{eq:prefix}.
Because $\log r^2<0$, the strict inequality $r^{2N}<c^2$ is equivalent
to $N>\log(c^2)/\log(r^2)$, which gives~\eqref{eq:N}.
The separated-sample result follows from Corollary~\ref{cor:separated}
with both observed indices retained in the prefix.
\end{proof}

The sharpness in Theorem~\ref{thm:two} concerns the observation constant.
The sufficient prefix in Theorem~\ref{thm:prefix} can be shortened for
specific coefficients or initial values.

\begin{theorem}[Adaptive rational certificate]\label{thm:adaptive}
Suppose $a,b,x_0,x_1$ are rational and satisfy~\eqref{eq:stable}.
After evaluating $x_0,\ldots,x_{n-1}$, define
$M_n^2=\max_{0\le j<n}x_j^2$. If
\begin{equation}\label{eq:stop}
 R^2r^{2n}<M_n^2,
\end{equation}
then $M_n$ is the infinite absolute peak. The all-zero trajectory is
recognized from $x_0=x_1=0$; every other trajectory eventually satisfies
\eqref{eq:stop}. Every operation and comparison in this certificate can be
performed with rational arithmetic.
\end{theorem}
\begin{proof}
For all future indices $j\ge n$,~\eqref{eq:envelope} and
\eqref{eq:stop} give $x_j^2<M_n^2$. Thus the observed peak is global.
A nonzero trajectory has a positive initial maximum, while $R^2r^{2n}$
tends to zero. Exact fractions represent the recurrence, squared radius,
powers of $r^2$, and the comparison, so no trigonometric or logarithmic
approximation is required by the algorithm.
\end{proof}

Strict comparison also certifies that all peak ties occur in the retained
prefix. It gives the earliest maximizing index by scanning the prefix in
order. The same certificate answers every finite-horizon query of length
at least $n$; shorter horizons use their observed prefix maximum.

\section{Optimal constants for a finite observation set}
Let $T\subset\Z_{\ge0}$ be a finite set and define row vectors
$\ell_k=(a_k,b_k)$. Assume these rows span $\R^2$.
Form the strip polygon
\begin{equation}\label{eq:polygon}
 \mathcal K_T=\{z\in\R^2:|\ell_kz|\le1\text{ for all }k\in T\}.
\end{equation}
It is bounded, contains a neighborhood of the origin, and is centrally
symmetric. Let $\operatorname{vert}(\mathcal K_T)$ be its vertices.

\begin{theorem}[Sharp finite-observation certificate]\label{thm:polygon}
The largest constant $c_T$ satisfying
\[
 \max_{k\in T}|x_k|\ge c_T R\quad\text{for every trajectory}
\]
is given by
\begin{equation}\label{eq:polygonconstant}
 \boxed{c_T^{-2}=\max_{z\in\operatorname{vert}(\mathcal K_T)}\norm z_2^2.}
\end{equation}
If $a,b$ are rational, $c_T^2$ is rational and can be obtained by checking
intersections of pairs of signed strip boundaries.
Any $N>\max T$ with $r^{2N}<c_T^2$ contains the infinite absolute peak.
\end{theorem}
\begin{proof}
For $z\ne0$, let $M(z)=\max_{k\in T}|\ell_kz|$.
The spanning assumption gives $M(z)>0$. Normalize by $M(z)$:
$z/M(z)\in\mathcal K_T$. Consequently,
\[
 \inf_{z\ne0}\frac{M(z)}{\norm z_2}
 =\left(\max_{w\in\mathcal K_T}\norm w_2\right)^{-1}.
\]
A convex function on a compact polytope attains a maximum at a vertex,
so the squared norm gives~\eqref{eq:polygonconstant}. Conversely a
maximizing vertex attains the bound, proving sharpness.
Each vertex is an intersection of two linearly independent active
boundaries $\ell_jz=\pm1$, $\ell_kz=\pm1$.
For rational coefficients the intersection and its squared norm are
rational. Enumerate such intersections and keep those satisfying every
strip. The final peak assertion repeats the proof of
Theorem~\ref{thm:prefix} with the observation maximum over $T$.
\end{proof}

Pairwise boundary enumeration uses $O(|T|^2)$ candidates and
$O(|T|)$ membership checks per candidate, giving a direct
$O(|T|^3)$ arithmetic implementation. This provides an exact elementary
construction; planar half-plane methods can exploit the geometry further.

\section{Finite and infinite quadratic sums}
The same two-dimensional state determines complete quadratic tails.
Let $Q$ be any symmetric real matrix. Since $r<1$, the convergent series
\begin{equation}\label{eq:stein}
 W=\sum_{n=0}^{\infty}(S^n)^\mathsf TQS^n
\end{equation}
is the unique solution of $W-S^\mathsf TWS=Q$.
For $K\ge0$, define $W_K=\sum_{n=0}^{K-1}(S^n)^\mathsf TQS^n$,
with $W_0=0$.

\begin{proposition}[Exact finite sum]\label{prop:energy}
\begin{align}
 W_K&=W-(S^K)^\mathsf TWS^K,\label{eq:finiteW}\\
 \sum_{n=0}^{K-1}z_n^\mathsf TQz_n&=z^\mathsf TW_Kz,
 \quad z_n=S^nz,\label{eq:finiteE}\\
 \sum_{n=0}^{K-1}z_n&=(I-S)^{-1}(I-S^K)z.\label{eq:linearsum}
\end{align}
Thus a cost $\sum_{n=0}^{K-1}(\kappa+\ell^\mathsf Tz_n+z_n^\mathsf TQz_n)$
is obtained from these expressions and $K\kappa$.
For rational input data, the result is rational. With fixed state dimension,
repeated squaring evaluates these formulas in $O(\log(K+1))$ arithmetic
operations after solving the fixed-size Stein system.
\end{proposition}
\begin{proof}
Multiplying~\eqref{eq:stein} on both sides by $S^K$ and its transpose
shifts the sum to indices $n\ge K$. Subtracting leaves~\eqref{eq:finiteW},
and substitution gives~\eqref{eq:finiteE}. The finite geometric-series
identity gives~\eqref{eq:linearsum}; $I-S$ is invertible because the
eigenvalues of $S$ have modulus below one. The Stein equation is a
nonsingular linear system with rational coefficients when the inputs
are rational. Binary exponentiation gives the arithmetic-operation bound.
\end{proof}

This is the classical Stein-sum identity, included to make the joint
peak-and-sum certificate self-contained. The arithmetic bound counts exact
field operations. Rational numerators and denominators can grow linearly
in $K$ in bit length, and the output representation must be charged in a
bit-complexity analysis.

\section{Worked certificate and verification}
Take $a=3/5$, $b=1/5$, $x_0=0$, $x_1=1$. Then
$r^2=2/5$, $R^2=25$ and $c^2=1/65$.
The uniform prefix bound gives $N=5$. The first four values are
\[
 0,\quad1,\quad\frac65,\quad\frac{26}{25}.
\]
Their peak is $6/5$. At the next index,
\[
 R^2r^8=25\left(\frac25\right)^4
 =\frac{16}{25}<\frac{36}{25}=M_4^2.
\]
Four rational samples therefore certify the infinite absolute peak,
first attained at index two. For $Q=e_1e_1^\mathsf T$, the Stein solution is
\[
 W=\begin{pmatrix}58/39&3/13\\3/13&7/39\end{pmatrix}.
\]
The initial state is $(0,5)^\mathsf T$, hence
$\sum_{n\ge0}x_n^2=175/39$.

The accompanying verifier checks 507 exact observation inequalities,
five sharp equality witnesses, 120 rational infinite-peak certificates
against 320-sample independent recurrence prefixes, and 720 exact finite
square sums against direct summation. A further 2420 rational-state checks
and twenty equality witnesses verify the strip-polygon construction. The infinite peak assertion in each
case is supplied by~\eqref{eq:stop}; the longer prefix is an independent
execution check. The source-derived figure uses the worked recurrence
and its exact rational envelope certificate.

\section{Context}
Finite determination of output constraints for stable linear systems is
part of the theory of maximal output-admissible sets; see
Gilbert and Tan~\cite{gilbert}. The present construction specializes to
conjugate second-order recurrences and gives a closed-form sharp
two-sample observation constant, its exact strip-polygon extension, and
a rational adaptive stopping certificate. The standard state representation
and the classical Stein relation~\cite{stein} supply the complete
finite-horizon cost alongside the peak.

\begin{thebibliography}{9}
\bibitem{gilbert} E. G. Gilbert and K. T. Tan,
Linear systems with state and control constraints: the theory and application
of maximal output admissible sets, IEEE Transactions on Automatic Control
36(9) (1991), 1008--1020.
\href{https://doi.org/10.1109/9.83532}{doi:10.1109/9.83532}.
\bibitem{stein} P. Stein, Some general theorems on iterants,
Journal of Research of the National Bureau of Standards 48(1)
(1952), 82--83. \href{https://doi.org/10.6028/jres.048.010}{doi:10.6028/jres.048.010}.
\end{thebibliography}
\end{document}
